\documentclass[main.tex]{subfiles}

\begin{document}
\section{Fialová, nikoli zelená planeta – Země mohla mít jinou barvu}
Oficiální barva Země, navzdory proklamacím politiků, je dnes kombinací zelené a~modré – to vychází ze způsobu, jakým světlo odrážejí dvě zásadní ingredience naší planety. Totiž voda v~různé podobě, a~rostliny obsahující chlorofyl. Máme tak v~hlavě jednoduchou rovnici fotosyntéza = chlorofyl = zelená barva.\par
Jenže tahle rovnice platí až pro relativně „moderní“ Zemi. Molekulární biolog Shiladitya DasSarma v~roce 2007 navrhl, že úplně první světlo-sklízející organismy mohly používat retinal – mnohem jednodušší pigment než chlorofyl. A ten světlo neabsorbuje stejně…
\subsection*{\enquote{Co se světlem dělá?}}
Retinal pohlcuje hlavně zeleno-žluté světlo, zatímco červené a~modré odráží. Barva je přitom dána tím, co nepohltí. Když tedy pigment sežere zelené fotony, zbyde směs červené a~modré – tedy fialová. To znamená, že kdyby jej dávní mikrobiočichové skutečně užívali, měly být fialový až purpurový vzhled buněk i~celých mikrobiálních povlaků. Pokud by takové organismy dominovaly povrchu planety, Země by z~dálky nevypadala zeleně, ale fialově!\par
Nejde o~čistou spekulaci. Retinal dnes používají například haloarchea – tedy mikroorganismy žijící v~extrémně slaných prostředích. Jejich protein bakteriorodopsin funguje jako světlem poháněná protonová pumpa – funguje bez kyslíku, bez přímé fixace uhlíku, ale extrémně efektivně. Je to jeden z~nejjednodušších bioenergetických systémů, jaké známe.\par
Možné stopy archeálních lipidů se přitom nacházejí i~ve velmi starých sedimentech z~archaika. To naznačuje, že organismy s~retinalem tu skutečně mohly být dřív než chlorofylová fotosyntéza – tedy ještě před Velkou kyslíkovou událostí, která planetu zaplavila kyslíkem a~zásadně proměnila globální ekosystémy, včetně masivního ústupu anaerobních organismů.
\subsection*{\enquote{Proč se to změnilo?}}
Jedna z~navržených možností je právě konkurence o~světlo. Pokud povrch planety nebo horní vrstvy vody pokryly mikrobiální povlaky využívající retinal, pohlcovaly hlavně zelené světlo a~fungovaly by jako spektrální filtr. Organismy žijící pod nimi by tak už zelené fotony téměř neviděly a~byly by vystaveny hlavně červenému a~modrému světlu, které retinal nevyužívá. Evoluční výhodu by proto získaly pigmenty schopné „jíst to, co zbylo“, místo aby soutěžily o~vyčerpaný zdroj.\par
Fialové mikroby tak nejprve „sežraly“ zelené světlo, pozdější bakterie se proto adaptovaly na zbylé červené a~modré vlnové délky – a~z toho se vyvinul chlorofyl. Jakmile se objevila fotosyntéza štěpící vodu, začal se hromadit kyslík, planeta prošla epizodami globálního ochlazení starý fialový svět skončil na okraji ekosystémů.\par
Ironií je, že chlorofyl se stal tak úspěšným, že se z~něj stala evoluční past. Je chemicky složitý, ale „vidí“ jen určitou část spektra – zelené světlo už nikdy zpátky nevyužil. Proto je dnešní biosféra zelená, i~když energeticky bohaté zelené fotony doslova odráží zpátky do vesmíru.
\subsection*{\enquote{A to je jistý?}}
Ne, je to  hypotéza – zajímavý scénář s~různou mírou svědčících argumentů, ale bez relativní jistoty teorií. Purple Earth hypothesis není dogma, ale provokativní nápad, který dává smysl chemicky, evolučně i~energeticky.\par
Btw. to není scénář jenom tak pro prdel – pokud hledáme život jen podle „zeleného signálu“ chlorofylu, můžeme minout celé třídy planet, kde život zvolil jednodušší cestu. Retinalové biosféry by z~dálky vypadaly červeno-modře, ne zeleně. Takže ufounci skutečně mohou být fialoví!\par
I kdyby dávná Země ale nebyla fialová, scénář dává logiku – a~mohl se odehrát někde jinde ve vesmíru. Jinými slovy to, co dnes považujeme za „normální“, je jen vítězná varianta v~dlouhé historii slepých evolučních pokusů…
\newpage
\end{document}